% BHU PYQ -- Manually transcribed & error-corrected
\documentclass[11pt,a4paper]{article}

\usepackage[a4paper,top=2.5cm,bottom=2.5cm,left=2.8cm,right=2.8cm,headheight=14pt]{geometry}
\usepackage{amsmath,amssymb}
\usepackage[T1]{fontenc}
\usepackage[utf8]{inputenc}
\usepackage{lmodern}
\usepackage{microtype}
\usepackage{fancyhdr}
\usepackage{enumitem}
\usepackage{hyperref}
\usepackage{lastpage}
\usepackage{parskip}

\hypersetup{colorlinks=true,urlcolor=black,linkcolor=black,
  pdftitle={BPT-504: Electronic Devices and Circuits -- BHU PYQ},pdfauthor={Banaras Hindu University}}

\pagestyle{fancy}\fancyhf{}
\renewcommand{\headrulewidth}{0.4pt}\renewcommand{\footrulewidth}{0.4pt}
\fancyhead[L]{\small   \textit{Banaras Hindu University}}
\fancyhead[R]{\small   \textit{BPT-504: Electronic Devices and Circuits}}
\fancyfoot[C]{\small Page  \thepage\ of \pageref{LastPage}}


\newlist{parts}{enumerate}{2}
\setlist[parts,1]{label=(\alph*),leftmargin=*,itemsep=5pt,topsep=3pt}
\setlist[parts,2]{label=(\roman*),leftmargin=*,itemsep=3pt,topsep=2pt}

\newcommand{\pts}[1]{\hfill   \textit{[#1]}}


\begin{document}


\begin{center}
  {\large\bfseries BANARAS HINDU UNIVERSITY}\\[2pt]
  {\normalsize B.Sc. (Hons.) Semester V Examination, 2016-17}\\[6pt]
  \rule{0.8\linewidth}{0.5pt}\\[6pt]
  {\large\bfseries Paper No. BPT-504: Electronic Devices and Circuits}\\[4pt]
  {\normalsize Time: 3 Hours\hfill Maximum Marks: 70}
\end{center}
\rule{\linewidth}{0.4pt}
\smallskip



\noindent   \textit{Instructions: Answer all questions. Symbols have their usual meanings.}
\bigskip




\noindent   \textbf{Question 1.} Answer in brief any   \textbf{five} of the following: \pts{5 $ \times$ 4 = 20}

\begin{parts}
  \item What is the significance of the threshold voltage ($V_T$) in an enhancement mode MOSFET? How is its magnitude reduced during fabrication?
  \item Show that negative feedback improves the stability of the voltage gain of an amplifier against parameter variations.
  \item Why is a Unijunction Transistor (UJT) not used as an amplifier? Explain the condition under which it is used as a relaxation oscillator.
  \item What is meant by Common Mode Rejection Ratio (CMRR) in the case of an OP-AMP? Why should its value be ideally infinite?
  \item Write the Boolean expression represented by the Karnaugh map (K-map) shown below and simplify it:
    \[
    
\begin{array}{c|c|c|c|c}
      A ackslash BC & 00 & 01 & 11 & 10 \\ \hline
      0 & 1 & 1 & 0 & 0 \\ \hline
      1 & 0 & 1 & 1 & 0
    \end{array}
    \]
  \item How is a Cathode Ray Oscilloscope (CRO) superior to ordinary measuring instruments? How does the Sync control help in obtaining a stationary pattern?
\end{parts}

\medskip\rule{\linewidth}{0.2pt}




\noindent   \textbf{Question 2.}

\begin{parts}
  \item With necessary theory and diagrams, explain the construction, working, and I-V characteristics of a Silicon Controlled Rectifier (SCR). Why can one not use germanium instead of silicon to construct a controlled rectifier? \pts{7}
  \item Explain the condition under which a UJT exhibits negative resistance. \pts{3.5}
  \item The firing voltage of a UJT is $15\, \text{V}$. It is connected to a series RC circuit to act as a relaxation oscillator with $R = 50\, \text{k}\Omega$ and $C = 2000\, \text{pF}$. If the source voltage is $40\, \text{V}$, determine the time period of the saw-tooth voltage across the capacitor $C$. \pts{3.5}
\end{parts}

\medskip\rule{\linewidth}{0.2pt}




\noindent   \textbf{Question 3.}

\begin{parts}
  \item Refer to the basic common-base BJT amplifier configuration. The h-parameter values for the transistor are:
    \[ h_{ib} = 40\,\Omega, \quad h_{rb} = 3  \times 10^{-4}, \quad h_{fb} = -0.99, \quad h_{ob} = 0.5\,\mu \text{mhos} \]
    Determine: (i) the closed-loop current gain taking source resistance $R_s = 600\,\Omega$ and load resistance $R_L = 10\, \text{k}\Omega$ into account, and (ii) the input resistance. \pts{6}
  \item Explain in brief the high-frequency hybrid-$\pi$ model of a CE transistor and discuss its advantages over the h-parameter model at high frequencies. \pts{8}
\end{parts}
\hfill   \textbf{OR}

\begin{parts}
  \item What do you mean by power amplifiers? Give the circuit diagram of class A, class B, and class C power amplifiers. Show that the maximum efficiency of a class B push-pull amplifier is $78.5\%$. \pts{8}
  \item What do you mean by non-linear distortion in a power amplifier circuit? Discuss a method to minimize such distortions. \pts{6}
\end{parts}

\medskip\rule{\linewidth}{0.2pt}




\noindent   \textbf{Question 4.}

\begin{parts}
  \item What do you mean by the term Amplitude Modulation (AM)? Derive the expression for an amplitude-modulated wave and discuss its frequency spectrum. \pts{8}
  \item Give the truth table of a full adder. Obtain the Boolean expressions for sum ($S$) and carry ($C_o$) from this table in SOP form. Simplify these expressions using a K-map and implement the circuit using basic logic gates. \pts{6}
\end{parts}

\medskip\rule{\linewidth}{0.2pt}




\noindent   \textbf{Question 5.}

\begin{parts}
  \item What is meant by feedback in amplifiers? Differentiate between positive and negative feedback. Draw the circuit diagram of a current-series feedback amplifier and show that negative feedback increases the input resistance and output resistance of the amplifier. \pts{8}
  \item How does the feedback network of a Hartley oscillator differ from the feedback network of a Colpitts oscillator? What will be the oscillator frequency in the two cases? \pts{6}
\end{parts}
\hfill   \textbf{OR}

\begin{parts}
  \item Explain the difference among astable, monostable, and bistable multivibrators in terms of their stable states and applications. \pts{7}
  \item Define the parameters relative amplification factor ($\mu$), drain resistance ($r_d$), and transconductance ($g_m$) in the case of a JFET amplifier, and establish the relationship among them. How is a JFET different from a bipolar junction transistor? \pts{7}
\end{parts}

\vfill
\rule{\linewidth}{0.4pt}

\begin{center}\small
  \href{https://bhu-exam.vercel.app/}{Click here to visit our website}\\[4pt]
  \href{https://me-aryan.vercel.app/}{Click here to visit our contributor}\\[4pt]
  \href{https://quashan-me.vercel.app/}{Click here to visit our contributor}
\end{center}

\end{document}
